\begin{thebibliography}{9}
\bibitem{aitchison} Aitchison, J. (1986). \textit{The Statistical Analysis of Compositional Data}. Chapman \& Hall.
\bibitem{martin} Martín-Fernández, J. A., Barceló-Vidal, C. y Pawlowsky-Glahn, V. (2003). Dealing with zeros and missing values in compositional data sets using nonparametric imputation. \textit{Mathematical Geology}, 35, 253--278. \url{https://doi.org/10.1023/A:1023866030544}.
\bibitem{egozcue} Egozcue, J. J., Pawlowsky-Glahn, V., Mateu-Figueras, G. y Barceló-Vidal, C. (2003). Isometric logratio transformations for compositional data analysis. \textit{Mathematical Geology}, 35, 279--300. \url{https://doi.org/10.1023/A:1023818214614}.
\bibitem{gap} Tibshirani, R., Walther, G. y Hastie, T. (2001). Estimating the number of clusters in a data set via the gap statistic. \textit{Journal of the Royal Statistical Society, Series B}, 63(2), 411--423. \url{https://doi.org/10.1111/1467-9868.00293}.
\bibitem{hennig} Hennig, C. (2007). Cluster-wise assessment of cluster stability. \textit{Computational Statistics \& Data Analysis}, 52(1), 258--271. \url{https://doi.org/10.1016/j.csda.2006.11.025}.
\bibitem{inei} INEI. \textit{Producto Bruto Interno por Departamentos}. Portal de estadísticas y cuentas departamentales, base 2007. \url{https://m.inei.gob.pe/estadisticas/indice-tematico/producto-bruto-interno-por-departamentos-9089/}. Consulta: 28 de octubre de 2026. Libros utilizados en \texttt{data/raw}; véase el manifiesto de fuentes.
\bibitem{linearmodels} Sheppard, K. y colaboradores. \textit{linearmodels: Panel models. Documentación de PanelOLS}. \url{https://bashtage.github.io/linearmodels/panel/model.html}. Versión de ejecución registrada en el manifiesto.
\end{thebibliography}
\clearpage
\appendix
\section{Clasificación y sensibilidades}
\begin{table}[htbp]\centering\small\begin{tabular}{p{.27\linewidth}p{.65\linewidth}}
\toprule
Perfil & Departamentos \\
\midrule
Diversificado/Agro & Lambayeque, Lima, San Martín \\
Servicios/Manufactura & Amazonas, Ayacucho, Cajamarca, Huánuco, Ica, La Libertad, Loreto, Piura, Puno, Tumbes, Ucayali \\
Minero/Diversificado & Arequipa, Cusco, Junín, Madre de Dios, Moquegua, Tacna, Áncash \\
Minero/Enclave & Apurímac, Huancavelica, Pasco \\
\bottomrule
\end{tabular}

\caption{Miembros de los perfiles Ward K=4}\end{table}
\begin{table}[htbp]\centering\small\begin{tabular}{rrrrr}
\toprule
K & Silueta Ward & Gap & EE Gap & Jaccard Ward \\
\midrule
1 & --- & 0.356 & 0.145 & --- \\
2 & 0.498 & 0.201 & 0.120 & 0.625 \\
3 & 0.350 & 0.140 & 0.148 & 0.736 \\
4 & 0.407 & 0.231 & 0.123 & 0.850 \\
5 & 0.367 & 0.057 & 0.164 & 0.786 \\
\bottomrule
\end{tabular}

\caption{Validación de Ward sobre PCA2}
\nota{Jaccard de submuestreo; no es una probabilidad de pertenencia. La regla Gap de un error estándar considera K=1.}\end{table}
\begin{table}[htbp]\centering\small\begin{tabular}{lr}
\toprule
Representación / reemplazo & ARI frente a PCA2 \\
\midrule
CLR+PCA2 & 1.000 \\
CLR+PCA3 & 0.567 \\
ILR-full & 0.867 \\
delta=0.01 & 1.000 \\
delta=0.05 & 1.000 \\
delta=0.1 & 1.000 \\
\bottomrule
\end{tabular}

\caption{Sensibilidad de K=4}
\nota{ARI frente a CLR+PCA2. Cambiar delta no perturba la composición porque los promedios observados no contienen ceros.}\end{table}
\FloatBarrier
\figura{fig3b_jerarquia_ward.pdf}{Jerarquía Ward y cortes derivados de las fusiones}{fig:jerarquia}{Colores de nombres según K=4. K=2 separa \WardSmallCount\ y \WardLargeCount\ unidades; no es la agregación económica \MacroMiningCount/\MacroOtherCount. Los cortes se calculan entre alturas consecutivas.}
\FloatBarrier
\section{Cargas de componentes principales}
\begin{table}[htbp]\centering\small\begin{tabular}{p{.58\linewidth}rrr}
\toprule
Actividad & PC1 & PC2 & PC3 \\
\midrule
Agricultura, Ganadería, Caza y Silvicultura & -0.130 & -0.542 & -0.373 \\
Extracción de Petróleo, Gas y Minerales & 0.906 & 0.137 & -0.241 \\
Manufactura & -0.152 & 0.630 & 0.212 \\
Electricidad, Gas y Agua & 0.145 & -0.289 & 0.825 \\
Construcción & 0.026 & -0.175 & 0.039 \\
Comercio & -0.181 & 0.016 & -0.216 \\
Transporte, Almacen., Correo y Mensajería & -0.110 & 0.152 & -0.073 \\
Alojamiento y Restaurantes & -0.090 & 0.169 & -0.023 \\
Telecom. y Otros Serv. de Información & -0.202 & 0.217 & -0.149 \\
Administración Pública y Defensa & -0.078 & -0.277 & 0.029 \\
Otros Actividades/1 & -0.133 & -0.038 & -0.029 \\
\bottomrule
\end{tabular}

\caption{Cargas PCA de las actividades armonizadas}
\nota{CLR centrado, sin estandarizar columnas. Los signos pueden invertirse sin alterar la geometría. Las cargas no son coeficientes de una regresión del VAB.}\end{table}
